%***************************************************************************************************
% preamble/floats-and-tables.tex - podpisy, rysunki, tabele.
%***************************************************************************************************

\graphicspath{{images/}{logos/}}

\captionsetup{font=small, labelfont=bf, labelsep=period, skip=5pt}
\captionsetup[figure]{position=bottom}
\captionsetup[table]{position=top}
\captionsetup[subfigure]{font=footnotesize, labelfont=bf, labelsep=period, skip=3pt}

\floatplacement{figure}{htbp}
\setlength{\intextsep}{8pt plus 2pt minus 2pt}
\setlength{\textfloatsep}{10pt plus 2pt minus 2pt}
\setlength{\floatsep}{8pt plus 2pt minus 2pt}

\renewcommand{\topfraction}{0.92}
\renewcommand{\bottomfraction}{0.85}
\setcounter{topnumber}{3}
\setcounter{bottomnumber}{3}
\setcounter{totalnumber}{4}
\renewcommand{\textfraction}{0.06}
\renewcommand{\floatpagefraction}{0.85}

\raggedbottom

\makeatletter
\setlength\@fptop{0\p@}
\makeatother

% ---------------------------------------------------------------------------------------------
% Rysunek: \Figure[szerokość]{plik}{podpis}{etykieta}
% ---------------------------------------------------------------------------------------------
\newcommand{\Figure}[4][1.0]{%
  \begin{figure}[H]
    \centering
    \IfFileExists{images/#2}{%
      \includegraphics[width=#1\linewidth, keepaspectratio]{images/#2}%
    }{%
      \PackageWarning{oss-fw-project}{Brak pliku obrazu images/#2}%
      \fbox{\parbox[c][5cm][c]{0.8\linewidth}{\centering\textit{Brak pliku obrazu}\\[1ex]\texttt{\detokenize{images/#2}}}}%
    }%
    \caption{#3}\label{#4}%
  \end{figure}%
}

% ---------------------------------------------------------------------------------------------
% Tabele
% ---------------------------------------------------------------------------------------------
\newcommand{\tablefont}{\footnotesize}
\renewcommand{\arraystretch}{1.08}
\setlength{\tabcolsep}{5pt}

\newcolumntype{L}[1]{>{\raggedright\arraybackslash}p{#1}}
\newcolumntype{C}[1]{>{\centering\arraybackslash}m{#1}}
\newcolumntype{R}[1]{>{\raggedleft\arraybackslash}p{#1}}
\newcolumntype{Y}{>{\raggedright\arraybackslash}X}
\newcolumntype{Z}{>{\centering\arraybackslash}X}
